\clearpage
\phantomsection% 
\addcontentsline{toc}{chapter}{ABSTRACT}
\thispagestyle{fancy}

\begin{center}
	\large \bfseries \MakeUppercase{Abstract}\\
	\normalsize \normalfont {\thetitleEN}\\
	\normalsize \normalfont {\theauthor}\\
	\bigskip
	
	\normalsize \normalfont \justifying \singlespacing
	Skin disease classification is one of the applications of artificial 
	intelligence in the medical field, particularly for automated 
	diagnosis through medical image analysis. The dataset used in this 
	study is DermaMNIST, which consists of seven classes. However, this 
	dataset has an issue related to class imbalance in the data 
	distribution, which may affect the performance of the model, 
	especially for minority classes. This condition causes the model to be 
	more accurate in predicting the majority classes. This study aims to 
	analyze the effect of applying data augmentation to minority classes 
	on the performance of deep learning models in skin disease image 
	classification. The CNN-based models used in this research are 
	ResNet-18 and ResNet-50, which were trained using two image sizes: 28x28 and 224x224. 
	The research process includes data 
	preprocessing, applying image augmentation to minority classes, model 
	training, and model performance evaluation using several metrics, 
	including accuracy, precision, recall, F1-score, AUC, and Matthews 
	Correlation Coefficient (MCC). The results show that the ResNet-18 
	model trained with 224x224 image size achieved an accuracy of 80\%, 
	precision macro 71\%, recall macro 76\%, F1-score macro 73\%, 
	AUC 95\%, and MCC 64\%. Based on these results, the application 
	of data augmentation on minority classes can reduce the impact of data 
	imbalance and improve the model's ability to classify skin disease images.

	% Halaman ABSTRACT berisi uraian tentang latar belakang, tujuan, metodologi penelitian, hasil / kesimpulan. Ditulis dalam BAHASA INGGRIS tidak lebih dari 250 kata, dengan jarak antar baris satu spasi. Secara khusus, kata dan kalimat pada halaman ini tidak perlu ditulis dengan huruf miring meskipun menggunakan Bahasa Inggris, kecuali terdapat huruf asing lain yang ditulis dengan huruf miring (misalnya huruf Latin atau Greek, dll). Pada akhir abstract ditulis kata “Keywords” yang dicetak tebal, diikuti tanda titik dua dan kata kunci yang tidak lebih dari 5 kata. Keywords terdiri dari kata-kata yang khusus menunjukkan dan berkaitan dengan bahan yang diteliti, metode/instrumen yang digunakan, topik penelitian. Keywords diketik pada jarak dua spasi dari baris akhir isi abstrak.\par
    
    \vspace{20pt}
	\raggedright\textbf{Keywords: Skin Disease Classification, Data Augmentation, ResNet-18, ResNet-50, DermaMNIST}
	
	\vfill
	
\end{center}
\clearpage